\begin{table}
    \centering
  % python examples/regen_gravity_quartet_tables.py (rsw_sphere.plotting.wave_set_table.wave_set_properties
    % + WaveSet.amplitudes_from_velocities, same generation discipline as Table \ref{tab: cap41})
    \begin{tabular}{|c|c|c|c|c|c|c|c|}\hline
    & Mode  & Freq. & Period (days) & Coeff.$_1$ &  Coeff.$_2$  & Zonal&$A_0$\\\hline
         a &  RH(4,5) &$-0.12524$&$3.992$& $0.3642$ & $-0.5580$ &30&$0.0918$  \\
         b &  RH(3,4) &$-0.13435$&$3.722$& $0.4440$ & $-0.5911$ &30&$0.0837$  \\
         c &  RH(1,2) &$-0.09944$&$5.028$& $-0.0287$ & $-$ &30&$0.0695$  \\
         d &  EG(1,1) &$0.36969$&$1.352$& $-$ & $-0.1326$ &0&$0.0000$  \\
        \hline
    \end{tabular}
    \caption{Four-wave configuration containing the Kelvin-like gravity mode EG(1,1) (Quartet C, Figure \ref{fig: topology_overview}c): each mode's frequency $\omega$, linear period (days), zonal velocity (m/s) and amplitude $A_0$ (eq. \ref{eq: Azonal}). Coeff.$_1$ is the RH-only triad (a,b,c); Coeff.$_2$ is the triad completed by EG(1,1) (a,b,d).}
    \label{tab: cap42}
\end{table}

\begin{table}
    \centering
  % python examples/regen_gravity_quartet_tables.py (rsw_sphere.plotting.wave_set_table.wave_set_properties
    % + WaveSet.amplitudes_from_velocities, same generation discipline as Table \ref{tab: cap41})
    \begin{tabular}{|c|c|c|c|c|c|c|c|}\hline
    & Mode  & Freq. & Period (days) & Coeff.$_1$ &  Coeff.$_2$  & Zonal&$A_0$\\\hline
         a &  RH(4,5) &$-0.12524$&$3.992$& $0.3642$ & $-0.7662$ &30&$0.0918$  \\
         b &  RH(3,4) &$-0.13435$&$3.722$& $0.4440$ & $-0.8999$ &30&$0.0837$  \\
         c &  RH(1,2) &$-0.09944$&$5.028$& $-0.0287$ & $-$ &30&$0.0695$  \\
         d &  EG(7,9) &$3.19551$&$0.156$& $-$ & $-0.3385$ &0&$0.0000$  \\
        \hline
    \end{tabular}
    \caption{Four-wave configuration containing the higher-frequency gravity mode EG(7,9) (Quartet D, Figure \ref{fig: topology_overview}d): columns as in Table \ref{tab: cap42}. Coeff.$_1$ is the RH-only triad (a,b,c); Coeff.$_2$ is the triad completed by EG(7,9), where EG(7,9) plays the SUM role ($m_d=m_a+m_b$), unlike Quartet C's EG(1,1).}
    \label{tab: cap43}
\end{table}
